\documentclass{article}
\usepackage{amsmath, amssymb, amsthm}
\usepackage{hyperref}
\usepackage{geometry}
\geometry{margin=1in}

\title{Erd\H{o}s Problem \#74}
\date{Last updated: 2025-08-31}
\author{Source: erdosproblems.com}

\begin{document}
\maketitle

\noindent\textbf{Status:} OPEN \quad \textbf{Prize: \$500}

\noindent\textbf{Tags:} graph theory, chromatic number, cycles

\medskip
\noindent\textbf{Problem Statement:}

Let $f(n)\to \infty$ (possibly very slowly). Is there a graph of infinite chromatic number such that every finite subgraph on $n$ vertices can be made bipartite by deleting at most $f(n)$ edges?

\bigskip
\noindent\textbf{Remarks:}

Conjectured by Erd\H{o}s, Hajnal, and Szemer\'{e}di \cite{EHS82}. 

R\"{o}dl \cite{Ro82} has proved this for hypergraphs, and also proved there is such a graph (with chromatic number $\aleph_0$) if $f(n)=\epsilon n$ for any fixed constant $\epsilon>0$.

It is open even for $f(n)=\sqrt{n}$. Erd\H{o}s offered \$500 for a proof but only \$250 for a counterexample. This fails (even with $f(n)\gg n$) if the graph has chromatic number $\aleph_1$ (see [111]).

\bigskip
\noindent\textbf{References:}

[EHS82] Erd\H{o}s, P. and Hajnal, A. and Szemer\'{e}di, E., On almost bipartite large chromatic graphs. Theory and practice of combinatorics (1982), 117-123.
 
 [Ro82] R\"{o}dl, Vojt\vEch, Nearly bipartite graphs with large chromatic number. Combinatorica (1982), 377-383.

\bigskip
\noindent\small{Source: \url{https://www.erdosproblems.com/74}}
\end{document}
